\documentclass[10pt,a4paper]{amsart}
\usepackage{amsmath,amssymb,amsthm,graphicx}
% Retained minimal preamble for the Stacks Project chapter flat.tex (More on Flatness).
% Verbatim excerpt of the chapter preamble: Xy-pic setup (native preamble lines 14-18) and the native
% theorem declarations with the mathematical macros (native preamble lines 154-217). Copied mechanically;
% nothing added, removed or redefined. The wrapper class and the amsmath/amssymb/amsthm/graphicx base
% packages are supplied by the assembler, so the class/externaldocument part of the native preamble is omitted.
% For commutative diagrams we use Xy-pic
\usepackage[all]{xy}

% We use 2cell for 2-commutative diagrams.
\xyoption{2cell}
% Theorem environments.
%
\theoremstyle{plain}
\newtheorem{theorem}[subsection]{Theorem}
\newtheorem{proposition}[subsection]{Proposition}
\newtheorem{lemma}[subsection]{Lemma}

\theoremstyle{definition}
\newtheorem{definition}[subsection]{Definition}
\newtheorem{example}[subsection]{Example}
\newtheorem{exercise}[subsection]{Exercise}
\newtheorem{situation}[subsection]{Situation}

\theoremstyle{remark}
\newtheorem{remark}[subsection]{Remark}
\newtheorem{remarks}[subsection]{Remarks}

\numberwithin{equation}{subsection}

% Macros
%
\def\lim{\mathop{\mathrm{lim}}\nolimits}
\def\colim{\mathop{\mathrm{colim}}\nolimits}
\def\Spec{\mathop{\mathrm{Spec}}}
\def\Hom{\mathop{\mathrm{Hom}}\nolimits}
\def\Ext{\mathop{\mathrm{Ext}}\nolimits}
\def\SheafHom{\mathop{\mathcal{H}\!\mathit{om}}\nolimits}
\def\SheafExt{\mathop{\mathcal{E}\!\mathit{xt}}\nolimits}
\def\Sch{\mathit{Sch}}
\def\Mor{\mathop{\mathrm{Mor}}\nolimits}
\def\Ob{\mathop{\mathrm{Ob}}\nolimits}
\def\Sh{\mathop{\mathit{Sh}}\nolimits}
\def\NL{\mathop{N\!L}\nolimits}
\def\CH{\mathop{\mathrm{CH}}\nolimits}
\def\proetale{{pro\text{-}\acute{e}tale}}
\def\etale{{\acute{e}tale}}
\def\QCoh{\mathit{QCoh}}
\def\Ker{\mathop{\mathrm{Ker}}}
\def\Im{\mathop{\mathrm{Im}}}
\def\Coker{\mathop{\mathrm{Coker}}}
\def\Coim{\mathop{\mathrm{Coim}}}

% Boxtimes
%
\DeclareMathSymbol{\boxtimes}{\mathbin}{AMSa}{"02}

%
% Macros for moduli stacks/spaces
%
\def\QCohstack{\mathcal{QC}\!\mathit{oh}}
\def\Cohstack{\mathcal{C}\!\mathit{oh}}
\def\Spacesstack{\mathcal{S}\!\mathit{paces}}
\def\Quotfunctor{\mathrm{Quot}}
\def\Hilbfunctor{\mathrm{Hilb}}
\def\Curvesstack{\mathcal{C}\!\mathit{urves}}
\def\Polarizedstack{\mathcal{P}\!\mathit{olarized}}
\def\Complexesstack{\mathcal{C}\!\mathit{omplexes}}
% \Pic is the operator that assigns to X its picard group, usage \Pic(X)
% \Picardstack_{X/B} denotes the Picard stack of X over B
% \Picardfunctor_{X/B} denotes the Picard functor of X over B
\def\Pic{\mathop{\mathrm{Pic}}\nolimits}
\def\Picardstack{\mathcal{P}\!\mathit{ic}}
\def\Picardfunctor{\mathrm{Pic}}
\def\Deformationcategory{\mathcal{D}\!\mathit{ef}}

% --- appended 2026-10-12 (Chief): the project own definitions of the reference/slogan/history
% display environments, copied mechanically from the shared Stacks preamble lines 8-12
% (\usepackage{verbatim} plus the three \newenvironment lines, which the chapter preamble keeps
% verbatim from the project preamble). Nothing added, removed or redefined: these environments
% carry no mathematical content and the project itself typesets them as comments. Without these
% three lines a counted statement whose own text uses them fails with Environment slogan undefined.
\usepackage{verbatim}
\newenvironment{reference}{\comment}{\endcomment}
\newenvironment{slogan}{\comment}{\endcomment}
\newenvironment{history}{\comment}{\endcomment}
\usepackage{cleveref}
\crefname{originallemma}{original lemma}{original lemmas}
\Crefname{originallemma}{Original lemma}{Original lemmas}
\makeatletter
\newlabel{algebra-lemma-extension}{{0519}{source}{}{sourcepointer-1}{}}
\newlabel{algebra-lemma-extension@cref}{{[originallemma][0][]0519}{[1][0][]source}{}{}{}}
\makeatother
\crefname{originallemma}{original lemma}{original lemmas}
\Crefname{originallemma}{Original lemma}{Original lemmas}
\makeatletter
\newlabel{algebra-lemma-unique-prime-over-localize-below}{{00EA}{source}{}{sourcepointer-2}{}}
\newlabel{algebra-lemma-unique-prime-over-localize-below@cref}{{[originallemma][0][]00EA}{[1][0][]source}{}{}{}}
\makeatother
\crefname{originaltheorem}{original theorem}{original theorems}
\Crefname{originaltheorem}{Original theorem}{Original theorems}
\makeatletter
\newlabel{algebra-theorem-projective-free-over-local-ring}{{0593}{source}{}{sourcepointer-3}{}}
\newlabel{algebra-theorem-projective-free-over-local-ring@cref}{{[originaltheorem][0][]0593}{[1][0][]source}{}{}{}}
\makeatother
\crefname{originallemma}{original lemma}{original lemmas}
\Crefname{originallemma}{Original lemma}{Original lemmas}
\makeatletter
\newlabel{lemma-fibres-irreducible-flat-projective-nonnoetherian}{{05FT}{source}{}{sourcepointer-4}{}}
\newlabel{lemma-fibres-irreducible-flat-projective-nonnoetherian@cref}{{[originallemma][0][]05FT}{[1][0][]source}{}{}{}}
\makeatother
\crefname{originallemma}{original lemma}{original lemmas}
\Crefname{originallemma}{Original lemma}{Original lemmas}
\makeatletter
\newlabel{lemma-induction-step}{{05I3}{source}{}{sourcepointer-5}{}}
\newlabel{lemma-induction-step@cref}{{[originallemma][0][]05I3}{[1][0][]source}{}{}{}}
\makeatother
\crefname{originalremark}{original remark}{original remarks}
\Crefname{originalremark}{Original remark}{Original remarks}
\makeatletter
\newlabel{remark-finite-presentation}{{05HY}{source}{}{sourcepointer-6}{}}
\newlabel{remark-finite-presentation@cref}{{[originalremark][0][]05HY}{[1][0][]source}{}{}{}}
\makeatother
\title{A module with a complete devissage that is flat at the point is finitely presented, free over the base and a direct sum of smooth pieces}
\author{The Stacks Project Authors}
\date{stacks-project:flat.tex:local-snapshot-20261011; GNU FDL 1.2 or later}
\begin{document}
\maketitle
\noindent\textit{Editorial note.} All delivered mathematical text is the author's own, copied from the pinned Stacks Project chapter file flat.tex; no statement or proof step was rewritten, reordered or removed.\par\medskip
\noindent\textit{Editorial note.} Preamble provenance: the chapter's own retained minimal preamble (sources/stacks-flat/retained-preamble.tex, lines 1-85) is a verbatim excerpt of the native Xy-pic setup and of the native theorem declarations and mathematical macros, followed by the project's own reference/slogan/history display-environment definitions copied mechanically from the shared preamble; the wrapper class and the amsmath/amssymb/amsthm/graphicx base packages are supplied by the assembler, and extra\_packages is left empty because everything needed is already in the retained preamble. No macro was invented or redefined.\par\medskip
\noindent\textit{Editorial note.} Same-chapter pointers are located at their \textbackslash{}label lines inside the pinned chapter file; cross-chapter pointers are located in the project's own tag list staged in the snapshot (sources/stacks-chow/tags.txt), at the verbatim lines carrying their canonical tags.\par\medskip
\noindent\textit{Allowed standard prerequisites (editorial).}\begin{itemize}
\item Finiteness of presentation along a complete devissage, and localisation at the unique prime lying over the point. Same chapter flat.tex, label remark-finite-presentation (proof line 2447) and Algebra, label algebra-lemma-unique-prime-over-localize-below (proof line 2452). The first identifies finiteness of presentation of N over S with that of M\_1 over B\_1, and the second identifies the localisations at the two primes, which together let (S, N) be replaced by (B\_1, M\_1).
\item The algebraic induction step for a flat map with geometrically integral fibres. Same chapter flat.tex, label lemma-induction-step, invoked at proof line 2456. Gives that the first map of the devissage is universally injective over the base and that the cokernel is flat at the point, which makes the induction hypothesis applicable to the cokernel.
\item Finiteness of presentation of an extension, freeness of a projective module over a local ring, and projectivity of the smooth pieces. Algebra, labels algebra-lemma-extension (proof line 2465) and algebra-theorem-projective-free-over-local-ring (proof line 2471); Same chapter flat.tex, label lemma-fibres-irreducible-flat-projective-nonnoetherian (proof line 2470). Shows that M\_1 is finitely presented as a B\_1-module once its cokernel is, that the smooth algebra B\_1 is free over the local base because it is projective, and that the decomposition of the cokernel supplies the required decomposition of M\_1.
\end{itemize}\medskip
\ifdefined\hypertarget\hypertarget{sourcepointer-1}{}\fi
\noindent\textit{Source reference.} Original lemma 0519: sources/stacks-chow/tags.txt, source lines 6527--6527. Navigation pointer for \textbackslash{}ref\{algebra-lemma-extension\} used inside the counted span (line 2465); the argument that consumes it is retained in the span.\par\medskip
\ifdefined\hypertarget\hypertarget{sourcepointer-2}{}\fi
\noindent\textit{Source reference.} Original lemma 00EA: sources/stacks-chow/tags.txt, source lines 527--527. Navigation pointer for \textbackslash{}ref\{algebra-lemma-unique-prime-over-localize-below\} used inside the counted span (line 2452); the argument that consumes it is retained in the span.\par\medskip
\ifdefined\hypertarget\hypertarget{sourcepointer-3}{}\fi
\noindent\textit{Source reference.} Original theorem 0593: sources/stacks-chow/tags.txt, source lines 6801--6801. Navigation pointer for \textbackslash{}ref\{algebra-theorem-projective-free-over-local-ring\} used inside the counted span (line 2472); the argument that consumes it is retained in the span.\par\medskip
\ifdefined\hypertarget\hypertarget{sourcepointer-4}{}\fi
\noindent\textit{Source reference.} Original lemma 05FT: sources/stacks-flat/pinned-original-source.tex, source lines 2247--2247. Navigation pointer for \textbackslash{}ref\{lemma-fibres-irreducible-flat-projective-nonnoetherian\} used inside the counted span (line 2470); the argument that consumes it is retained in the span.\par\medskip
\ifdefined\hypertarget\hypertarget{sourcepointer-5}{}\fi
\noindent\textit{Source reference.} Original lemma 05I3: sources/stacks-flat/pinned-original-source.tex, source lines 2356--2356. Navigation pointer for \textbackslash{}ref\{lemma-induction-step\} used inside the counted span (line 2456); the argument that consumes it is retained in the span.\par\medskip
\ifdefined\hypertarget\hypertarget{sourcepointer-6}{}\fi
\noindent\textit{Source reference.} Original remark 05HY: sources/stacks-flat/pinned-original-source.tex, source lines 1685--1685. Navigation pointer for \textbackslash{}ref\{remark-finite-presentation\} used inside the counted span (line 2447); the argument that consumes it is retained in the span.\par\medskip
\begin{lemma}
\label{lemma-complete-devissage-flat-finite-type-module}
Let $(R, \mathfrak m)$ be a local ring.
Let $R \to S$ be a ring map of finite presentation.
Let $N$ be a finite $S$-module.
Let $\mathfrak q$ be a prime of $S$ lying over $\mathfrak m$.
Assume that $N_{\mathfrak q}$ is flat over $R$, and
assume there exists a complete d\'evissage of $N/S/R$ at $\mathfrak q$.
Then $N$ is a finitely presented $S$-module, free as an $R$-module,
and there exists an isomorphism
$$
N \cong B_1^{\oplus r_1} \oplus \ldots \oplus B_n^{\oplus r_n}
$$
as $R$-modules where each $B_i$ is a smooth $R$-algebra with geometrically
irreducible fibres.
\end{lemma}

\begin{proof}
Let $(A_i, B_i, M_i, \alpha_i, \mathfrak q_i)_{i = 1, \ldots, n}$
be the given complete d\'evissage. We prove the lemma by induction on $n$.
Note that $N$ is finitely presented as an $S$-module if and only if
$M_1$ is finitely presented as an $B_1$-module, see
Remark \ref{remark-finite-presentation}.
Note that $N_{\mathfrak q} \cong (M_1)_{\mathfrak q_1}$ as $R$-modules
because (a) $N_{\mathfrak q} \cong (M_1)_{\mathfrak q'_1}$ where
$\mathfrak q'_1$ is the unique prime in $A_1$ lying over $\mathfrak q_1$
and (b) $(A_1)_{\mathfrak q'_1} = (A_1)_{\mathfrak q_1}$ by
Algebra, Lemma \ref{algebra-lemma-unique-prime-over-localize-below},
so (c) $(M_1)_{\mathfrak q'_1} \cong (M_1)_{\mathfrak q_1}$.
Hence $(M_1)_{\mathfrak q_1}$ is a flat $R$-module. Thus we may replace
$(S, N)$ by $(B_1, M_1)$ in order to prove the lemma. By
Lemma \ref{lemma-induction-step}
the map $\alpha_1 : B_1^{\oplus r_1} \to M_1$ is $R$-universally injective
and $\Coker(\alpha_1)_{\mathfrak q}$ is $R$-flat.
Note that $(A_i, B_i, M_i, \alpha_i, \mathfrak q_i)_{i = 2, \ldots, n}$
is a complete d\'evissage of $\Coker(\alpha_1)/B_1/R$ at
$\mathfrak q_1$. Hence the induction hypothesis
implies that $\Coker(\alpha_1)$ is finitely presented as a
$B_1$-module, free as an $R$-module, and has a decomposition as in the lemma.
This implies that $M_1$ is finitely presented as a $B_1$-module, see
Algebra, Lemma \ref{algebra-lemma-extension}.
It further implies that
$M_1 \cong B_1^{\oplus r_1} \oplus \Coker(\alpha_1)$
as $R$-modules, hence a decomposition as in the lemma.
Finally, $B_1$ is projective as an $R$-module by
Lemma \ref{lemma-fibres-irreducible-flat-projective-nonnoetherian}
hence free as an $R$-module by
Algebra, Theorem \ref{algebra-theorem-projective-free-over-local-ring}.
This finishes the proof.
\end{proof}

\end{document}
